\documentclass[11pt]{article}

\usepackage{amsmath,amssymb,amstext,amsthm}
\usepackage{geometry}
\usepackage{fancyhdr}
\usepackage[pdftex]{graphicx}
\usepackage{xcolor}

\geometry{
  verbose,
  dvips,
  width=430pt, marginparsep=0pt, marginparwidth=0pt,
  top=90pt, headheight=12pt, headsep=20pt, footskip=30pt, bottom=80pt}

\setlength{\parskip}{\medskipamount}
\setlength{\parindent}{0pt}

\linespread{1.05}

\newtheorem{theorem}{Theorem}
\newtheorem{lemma}[theorem]{Lemma}
\newtheorem{proposition}[theorem]{Proposition}
\newtheorem{corollary}[theorem]{Corollary}
\newtheorem{conjecture}[theorem]{Conjecture}
\newtheorem{obs}{Observation}
\newtheorem{clm}{Claim}
\newtheorem{ques}{Question}
\newtheorem{problem}{Problem}

\newtheorem{ex}{Example}


\newcommand{\spc}{\hspace{0.08in}}
\newcommand{\vs}{\vspace*{.1in}}
\newcommand{\E}{\mathbb{E}}
\newcommand{\Prob}{\mathbb{P}}
\newcommand{\edit}[1]{\emph{\textcolor{blue}{#1}}}
\newcommand*\dif{\mathop{}\!\mathrm{d}}


\renewenvironment{proof}{{\noindent \textbf{\textit{Proof.}}}}
{\hfill $\Box$ \vspace*{0.1in}}
\newenvironment{proofc}{{\noindent \textit{Proof of Claim.}}}
{\hfill $\Box$}



\begin{document}

\begin{LARGE}\begin{center}\bf 16.4 Green's Theorem\end{center}
\end{LARGE}

\vs\vs



\section{Warm Up}


\section{Motivation}
In the last section we looked at what happens when we go to take a line integral over a vector field and my vector field is conservative. This took care of a special case of the line integral problem and with the FToC we solved it! What if instead of taking a special vector field and any curve, we reversed that and looked at any vector field and a specific curve? Can we also answer the line integral problem?



\section{Definitions \& Theorems}

\begin{enumerate}
\item A closed simple curve is one that starts and ends at the same point and does not cross over itself. We also assume we travel in the counter-clockwise direction.

\item Green's Theorem says the following

\begin{theorem}(Green's Theorem)
If $C$ is a simple closed curve containing the region $R$ then,

\begin{equation*}
       \oint_C F \cdot \dif{r} =  \oint_C P \dif{x} + Q \dif{y} = \iint_R \left[\frac{\partial Q}{\partial x} - \frac{\partial P}{\partial y}\right] \dif{A}.
\end{equation*}
\end{theorem}




\item Just to make sure we see the various forms of a line integral and our notation is consistent, let $F(x,y) = P\hat{i} + Q \hat{j}$ and let our curve be defined by $\vec{r}(t) = x\hat{i} + y\hat{j}$. Remember $\vec{r'}(t) = f_x \hat{i} + f_y \hat{j}$. Then $F\cdot \vec{r'} = P f_x + Q f_y$


\item Okay what if we had the vector field $F(x,y) = -y \hat{i} + x\hat{j}$. For now this just seems like a regular vector field. But now consider the line integral through through a closed loop over this vector field! (Well half of it)

\begin{equation*}
    \frac{1}{2}\oint_C -y \dif{x} + x\dif{y} = \iint_R \left(  \frac{\partial Q}{\partial x}    - \frac{\partial P}{\partial y} \right) \dif{A} = \frac{1}{2} \iint_R 1 - (-1) \dif{A} = \iint_R 1 \dif{A}.
\end{equation*}

We now have an equation for the area bounded by a closed simple curve in terms of line integrals! This is cool!





\end{enumerate}




\section{Examples}


\begin{problem}
    Evaluate the following line integral where $C$ is the closed curve around the triangle with vertices $(0,0),(1,1), (0,1)$.

    \begin{equation*}
        \oint_C x^3 \dif{x} + xy \dif{y}
    \end{equation*}
\end{problem}


\vspace{7cm}

\begin{problem}
    Evaluate the following line integral where $C$ is the closed curve around $y=x$ and $y=x^2$.

    \begin{equation*}
        \oint_C (x^2y+x^3) \dif{x} + 2xy \dif{y}
    \end{equation*}
\end{problem}


\vspace{8cm}

\begin{problem}
    Evaluate the following line integral.

    \begin{equation*}
        \oint_C xe^x \dif{x} + y^5 \dif{y}
    \end{equation*}
\end{problem}


\vspace{6cm}



\begin{problem}
    Evaluate the following line integral where $C$ is the closed curve around the top half semicircle $x^2+y^2=1$ and the $x$-axis.

    \begin{equation*}
        \oint_C x^2y \dif{x} + y^3 \dif{y}
    \end{equation*}
\end{problem}

\vspace{10cm}

\begin{problem}
    Evaluate the following line integral where $C$ is the closed curve around the circle with radius 2.

    \begin{equation*}
        \oint_C 6xy \dif{x} + (3x^2+\ln{(1+y)}) \dif{y}
    \end{equation*}
\end{problem}

\vspace{9cm}

\begin{problem}
    Find the area of the ellipse $\frac{x^2}{a^2}+ \frac{y^2}{b^2} = 1$
\end{problem}


\vspace{9cm}

\begin{problem}
    Find the work done through the vector field $F(x,y) = \frac{y}{1+x^2}\hat{i} + (x+\tan^{-1}(x)) \hat{j}$ along the right loop of the curve $r^2 = \cos(2\theta)$.
\end{problem}


\newpage

\section{Scratch Work}

\end{document} 